\section{Final Report}

The final report is the main written record of your project. It should present the full product development process in a self-contained way so that a reader can understand the problem, the evidence gathered, the design decisions made, and the proposed final concept without needing to read your weekly submissions first.

Submit one group report in PDF format through ITEL. The cover page should identify the group members, project title, and any other details requested by the lecturer. The final report should normally be about 20 to 25 pages long, excluding appendices, unless the lecturer provides a different limit.

\subsection{Recommended Format and Tools}

You may use \LaTeX{} or another suitable word processor. \LaTeX{} is recommended because it supports consistent formatting, references, tables, and figure management, but it is not compulsory if your group can produce an equally professional report in another format.

Regardless of the software used, the report should be:
\begin{itemize}
    \item clearly structured with section headings;
    \item written in formal academic English;
    \item supported by captions, references, and labelled figures or tables; and
    \item checked carefully for consistency, grammar, and technical accuracy.
\end{itemize}

\subsection{Suggested Report Structure}

The structure below is recommended because it aligns with the weekly APID workflow and the common grading criteria.

\begin{enumerate}
    \item \textbf{Title and Executive Summary}
    \begin{itemize}
        \item State the product title clearly.
        \item Summarise the problem, proposed solution, and key result in a short opening section.
    \end{itemize}

    \item \textbf{Introduction and Problem Context}
    \begin{itemize}
        \item Explain the hydroponics-related problem or opportunity.
        \item Introduce the target users and why the project matters.
        \item Preview the report structure.
    \end{itemize}

    \item \textbf{Opportunity Identification and Product Planning}
    \begin{itemize}
        \item Present the innovation charter, identified opportunities, and planning rationale.
        \item Define the mission statement, target market, constraints, and stakeholders.
    \end{itemize}

    \item \textbf{Customer Needs and Product Specifications}
    \begin{itemize}
        \item Show how customer statements were translated into interpreted needs.
        \item Present measurable specifications, metric justifications, and benchmarking outputs.
    \end{itemize}

    \item \textbf{Concept Development}
    \begin{itemize}
        \item Describe functional decomposition, solution search, concept generation, and concept selection.
        \item Explain why the preferred concept was selected.
    \end{itemize}

    \item \textbf{Concept Testing and Prototyping Plan}
    \begin{itemize}
        \item Explain how the concept would be validated.
        \item Outline the prototype objective, scope, and next development steps.
    \end{itemize}

    \item \textbf{Intellectual Property and Conclusion}
    \begin{itemize}
        \item Comment on patent relevance or other ownership issues.
        \item Conclude with the main design outcome and lessons learned.
    \end{itemize}
\end{enumerate}

\subsection{Expected Evidence}

The final report should integrate the strongest outputs from the weekly sessions. Typical evidence includes:
\begin{itemize}
    \item interview or survey findings translated into customer needs;
    \item a needs-importance table and a needs-metric table with justified units;
    \item a needs-metric matrix and competitor benchmarking summary;
    \item function diagrams, solution tables, and concept combination outputs; and
    \item sketches, annotated diagrams, or prototype planning visuals.
\end{itemize}

Table~\ref{tab:session0-final-report-needs} shows a simplified example of how customer needs can be documented.

\begin{table}[htbp]
\centering
\caption{Example format for translating customer statements into interpreted needs}
\label{tab:session0-final-report-needs}
\renewcommand{\arraystretch}{1.2}
\begin{tabularx}{\textwidth}{|>{\raggedright\arraybackslash}p{0.24\textwidth}|>{\raggedright\arraybackslash}X|>{\raggedright\arraybackslash}X|}
\hline
\textbf{Prompt or theme} & \textbf{Customer statement} & \textbf{Interpreted need} \\
\hline
Typical use & ``I need to check the nutrient tank too often.'' & The system should make nutrient-level monitoring quick and easy. \\
\hline
Dislike & ``Water flow becomes inconsistent when the pump clogs.'' & The system should maintain stable nutrient flow and allow easy maintenance. \\
\hline
Suggested improvement & ``I want to know immediately when the pH changes.'' & The system should provide timely pH monitoring and alert functions. \\
\hline
\end{tabularx}
\end{table}

Table~\ref{tab:session0-final-report-metrics} gives a matching example for the specification stage.

\begin{table}[htbp]
\centering
\caption{Example format for linking needs to measurable metrics}
\label{tab:session0-final-report-metrics}
\renewcommand{\arraystretch}{1.2}
\begin{tabularx}{\textwidth}{|>{\raggedright\arraybackslash}X|>{\raggedright\arraybackslash}p{0.2\textwidth}|>{\raggedright\arraybackslash}p{0.12\textwidth}|>{\raggedright\arraybackslash}X|}
\hline
\textbf{Customer need} & \textbf{Metric} & \textbf{Unit} & \textbf{Justification} \\
\hline
Stable nutrient delivery & Nutrient flow-rate variation & L/min & Indicates whether the delivery system remains consistent during operation. \\
\hline
Easy monitoring & Display response time after sensor update & s & Shows how quickly a user receives useful system feedback. \\
\hline
Safe operation & Leakage rate during 24-hour operation & mL/h & Relates directly to reliability, cleanliness, and maintenance burden. \\
\hline
\end{tabularx}
\end{table}

\subsection{Submission Requirements}

Before submission, confirm that the report:
\begin{itemize}
    \item is complete and internally consistent;
    \item uses captions and references correctly;
    \item cites external sources where needed;
    \item follows any naming convention announced on ITEL; and
    \item reflects your own checked and edited understanding, including when LLM tools were used.
\end{itemize}

\subsection{About the Use of LLMs in the Report}

LLMs may help you improve structure, grammar, or brainstorming efficiency, but they cannot replace technical judgement. If you use an LLM while preparing the final report, verify every claim, remove unsupported content, and rewrite the material so that it matches your project evidence and the actual task requirements. Weak reports often result from unedited LLM output that sounds polished but does not answer the question accurately.
